\documentclass[../DoAn.tex]{subfiles}
\begin{document}

\begin{center}
    \Large{\textbf{ABSTRACT}}\\
\end{center}
\vspace{1cm}
Smartphone security today relies mainly on one-time authentication at unlock, such as PIN, fingerprint, or face recognition, so a device can still be accessed by an unauthorized person after it has been unlocked. Continuous authentication based on behavioral biometrics has been proposed to address this limitation; systems such as IntelliAuth, B2auth, M2Auth, and MultiLock exploit inertial sensors or touch interaction, but they remain limited in adapting to hardware differences across devices and lack a fallback authentication mechanism to replace the PIN when a stranger is detected.

This thesis adopts a multimodal approach that combines inertial sensors (\\
accelerometer, gyroscope, magnetometer) with touchscreen interaction (tap, scroll, keystroke), and exploits activity context by collecting activity-labeled data for three activities and then training and evaluating per context.

The proposed solution is a complete pipeline from data collection to deployment. The main branch, based on inertial sensors, uses an encoder that learns an embedding space; each owner is enrolled with a few representative embeddings (anchors) and is authenticated by cosine similarity, operating under an open-set recognition scenario. Confidence scores are smoothed over time using EWMA, and a fallback authentication mechanism based on a personal phone-shake gesture is triggered when the score falls below a threshold. The system is implemented in two Android applications together with a training pipeline in Python and TensorFlow.

The main contribution is a continuous authentication system operating in an open-set scenario together with a fallback mechanism, evaluated on a real dataset of 26 participants comprising 136{,}746 inertial windows, 34{,}836 touch events, and 22{,}694 keystroke events. Three architectures, CNN 1D, CNN-LSTM, and CNN-BiLSTM, are compared, and the decision functions that convert embeddings into confidence scores are studied systematically. The CNN 1D encoder using only inertial sensors is selected for deployment; with a cohort-based score-normalization decision function, this configuration attains an equal error rate of about 0.3\% in the closed-set scenario and about 1.4\% in the open-set scenario under leave-users-out cross-validation, a substantial improvement over the baseline decision function and without retraining the model. Adding the touch branch through linear score-level fusion was evaluated as a control but did not improve open-set performance, so the final deployment uses only the inertial branch.

\end{document}